\chapter{Legacy Sources}\label{app:legacy}
Earlier files are preserved unmodified in the project.

\begin{itemize}
\item \texttt{toarabic.ahk}, an AutoHotkey layout from the author's repository. It assigns one Arabic letter to each Latin key, so \texttt{p} and \texttt{v} map to key assignments and not to claims about how the sounds are written.
\item \texttt{pronouncing\_translit\_chat\_version.py}, the pronouncing-based script as pasted from a chat. Its defects are that the vowel table mixes letters and diacritics, the entry for \texttt{sh} is never reached, the voiced-$v$ test indexes letters with a phoneme index, and no vowel or sukoon placement is done.
\item \texttt{cli\_main\_as\_pasted.py}, the command-line wrapper, which imports \texttt{transliterate\_text} from a module named \texttt{transliterator}. A shim of that name forwards to the package.
\end{itemize}

\section{The old version}
The author's earlier program (the files \texttt{old\_version\_}\allowbreak\texttt{transliterator.py} and \texttt{old\_version\_}\allowbreak\texttt{mappings.py}) converts spelling to the first CMU pronunciation and then writes one letter and one mark per phoneme, with nine fixed words. Its consonant choices are $p\to$\ar{ب}, $v\to$\ar{ف}, hard $g\to$\ar{ج}, \emph{ch}$\to$\ar{ش}, \emph{ng}$\to$\ar{ن} and \emph{zh}$\to$\ar{ز}. It merges $p/b$, $v/f$, hard $g$ and $j$, \emph{ch} and \emph{sh}, and \emph{ng} and $n$. It writes each vowel letter before its mark (so AY gives \ar{يَ}), uses no sukoon outside the fixed words, and takes the first dictionary pronunciation without regard to context. Unrecognised phonemes are dropped silently. The author describes it as not definitive.

The current toolkit differs in three ways that matter. It places the vowel mark on the preceding consonant and writes glides with a sukoon (\ar{سَايْ}), which is the form on the handwritten page. It writes hard $g$ as \ar{غ} and \emph{ch} as \ar{تْش} by default, with switches (\texttt{--g jeem}) to return to the old choices. And the $v$ choice is a switch, with \ar{ف} the default as in the old version, since the author regards $v$ as a voiced $f$; \emph{individual} and \emph{universal} in the typed passage use \ar{ب} and are reproduced by \texttt{--v b}. Both versions use only standard Arabic letters, and the author accepts near equivalents such as \ar{ب} for $p$ when exact recovery of short words is not possible.
